%\ifodd\value{page}\hbox{}\newpage\fi
\chapter*{Ṛgveda 10.125 — Devī Sūkta, Mantra 2}
\addcontentsline{toc}{chapter}{Devī Sūkta 10.125.2}

\begin{sloka}
{\sanskritfont
\begin{tabular}{c}
अहं राष्ट्र्री सङ्गमनी वसूनां चिकित्वेषी प्रथमायज्ञियानाम् ।\\
तां मा देवा व्यदधुः पुरुत्रा भूूरिस्थात्रां भूूर्यावेशयन्तीम् ॥२॥
\end{tabular}

\vspace{1.5em}

{\middlesize \iastfont
ahaṁ rāṣṭrī saṅgamanī vasūnāṁ cikitveṣī prathamā yajñiyānām |\\
tāṁ mā devā vyadadhuḥ purutrā bhūri-sthātrāṁ bhūry āveśayantīm ||2||}

\vspace{1em}

{\middlesize
\anvaya{%
अहं राष्ट्र्री ।
अहं वसूनां सङ्गमनी ।
अहं प्रथमायज्ञियानां चिकित्वेषी ।
देवाः मां पुरुत्रा व्यदधुः ।
भूरिस्थात्रां भूरि आवेशयन्तीम् ।
}
}

\middlesize \iastfont \itmean{%
«Io sono la Sovrana, colei che riunisce i beni;  
colei che conosce e ricerca, la prima tra coloro che sono degni del sacrificio.  
Gli Dei mi hanno posta in molti luoghi,  
rendendomi vasta nel dimorare e capace di pervadere ampiamente.»%
}

}
\end{sloka}

\wordmeaning{%
\wordline{अहम् राष्ट्र्री}
{aham rāṣṭrī}
{ io sono la Sovrana (\textit{rāṣṭrī}), colei che governa e regge l’ordine; }%

\wordline{वसूनां सङ्गमनी}
{vasūnāṁ saṅgamanī}
{ colei che riunisce (\textit{saṅgamanī}) i beni e le ricchezze (\textit{vasu}), in senso materiale e cosmico; }%

\wordline{चिकित्वेषी प्रथमायज्ञियानाम्}
{cikitveṣī prathamā yajñiyānām}
{ colei che conosce e cerca la conoscenza (\textit{cikitveṣī}), la prima tra coloro che sono degni del sacrificio (\textit{yajñiya}); }%

\wordline{देवा व्यदधुः पुरुत्रा}
{devā vyadadhuḥ purutrā}
{ gli Dei (\textit{devāḥ}) mi hanno collocata (\textit{vyadadhuḥ}) in molti luoghi (\textit{purutrā}); }%

\wordline{भूरिस्थात्रां भूूर्यावेशयन्तीम्}
{bhūri-sthātrāṁ bhūry āveśayantīm}
{ rendendomi ampia nel dimorare (\textit{bhūri-sthātrā}) e capace di pervadere ampiamente (\textit{bhūry āveśayantīm}); }%
}

Nel secondo mantra, la voce della Dea assume un profilo apertamente regale e ordinatore. Definirsi \textit{rāṣṭrī} significa affermarsi come principio di sovranità, non politica ma cosmica: colei che regge e coordina. Essere “colei che riunisce i \textit{vasu}” non indica un semplice possesso di beni, ma la capacità di raccogliere, tenere insieme e rendere disponibili le potenze del mondo. La Dea è anche detta \textit{cikitveṣī}: colei che conosce e che cerca la conoscenza, posta al primo posto tra ciò che è degno del sacrificio. La conoscenza e il rito non sono separati, ma convergono nella sua presenza. Quando gli Dei “la pongono in molti luoghi”, non la distribuiscono come un oggetto: riconoscono piuttosto la sua natura pervasiva. La Dea dimora ampiamente e pervade ampiamente; è stabile e al tempo stesso diffusiva. Questo mantra chiarisce così che la sua sovranità non è localizzata: essa si manifesta ovunque come principio di coesione, conoscenza e ordine.

Il termine \textit{cikitveṣī} è uno dei più densi e significativi dell’intero \textit{Devī Sūkta}. Deriva dalla radice verbale \textit{cit / cikit-} (“percepire, comprendere, essere consapevoli”) unita al suffisso desiderativo \textit{-eṣī}, e indica colei che conosce e, nello stesso tempo, cerca la conoscenza. Non si tratta dunque di un sapere statico o posseduto una volta per tutte, ma di una intelligenza viva, in atto, che continuamente si volge verso il comprendere. Applicato alla Dea, \textit{cikitveṣī} esprime una forma di coscienza che è al tempo stesso fondamento e movimento: la conoscenza non come accumulo di contenuti, ma come principio ordinatore che interroga, discerne e rende il cosmo intelligibile. \textit{Cikitveṣī} definisce così una sovranità cognitiva, una regalità della coscienza, in cui il sapere è inseparabile dalla responsabilità cosmica.
